\section{Rules}

\newcommand{\universe}{\mathcal{U}}
\newcommand{\type}{\ \mathrm{type}}
\def\judg#1:#2{#1 \oftype #2}


\[
\begin{array}{l}
    \inferrule*[right=\text{inl-intro}]{
        \Gamma \vdash A, B \type}{
        \Gamma \vdash \mathrm{inl}(A, B) \oftype A \mid B} \\
    \\
    \inferrule*[right=\text{inr-intro}]{
        \Gamma \vdash A, B \type}{
        \Gamma \vdash \mathrm{inr}(A, B) \oftype A \mid B} \\
    \\
    \inferrule*[right=\text{inl-elim}]{
        \Gamma \vdash A, B \type \quad \Gamma \vdash x = \mathrm{inl}(A, B)}{
        \Gamma \vdash \mathrm{unpack}(x) \oftype A} \\
    \\
    \inferrule*[right=\text{inr-elim}]{
        \Gamma \vdash A, B \type \quad \Gamma \vdash x = \mathrm{inr}(A, B)}{
        \Gamma \vdash \mathrm{unpack}(x) \oftype B} \\
    \\
    \inferrule*[right=\text{pair-intro}]{
        \Gamma \vdash A, B \type}{
        \Gamma \vdash \mathrm{pair(A, B)} : A \times B} \\
    \\
    \inferrule*[right=\text{$\pi_1$-elim}]{
        \Gamma \vdash A, B \type}{
        \Gamma \vdash \pi_1(\mathrm{pair(A, B)}) \oftype A} \\
    \\
    \inferrule*[right=\text{$\pi_2$-elim}]{
        \Gamma \vdash A, B \type}{
        \Gamma \vdash \pi_2(\mathrm{pair(A, B)}) \oftype B} \\
    \\
    \inferrule*[right=\text{lambda-intro}]{
        \Gamma \vdash A, B \type \quad \Gamma, x:A \vdash M:B}{
        \Gamma \vdash \lambda x:A.M \oftype \Pi x:A.B
    } \\
    \\
    \inferrule*[right=\text{lambda-elim}]{
        \Gamma \vdash A, B \type \quad \Gamma \vdash M:\Pi x:A.B \quad \Gamma \vdash N:A}{
        MN \oftype B[x := N]
    }\\
\end{array}
\]

Elimination rules are explicative and introduction rules are used to create hypothetical proofs of an existing type or a new type.

anytime algorithms